% Editorial proposal only. Keep the existing scientific figure/table payloads
% unchanged at the indicated insertion points. No canonical file was edited.
\section{Experiments}
Useful graph adaptation requires three things: the simulator must be able to use additional messages, the controller must place them where they improve a prediction, and that improvement must persist as predictions feed back into the next state. We test these requirements in turn. The first comparison holds the expansion policy fixed and changes training. The second freezes each trained model and compares policies on identical observed histories. Full rollouts then expose the controller to the states it creates.

\subsection{Can the simulator learn to use extra messages?}
We pair base-only training with mixed-graph training, which expands a randomly chosen half of training examples by adding a uniformly sampled quarter of the optional annulus pairs. Both arms use the faithful objective and matched initialization, frames and noise. Goop and Sand train from initialization for 100k updates; WaterDrop uses paired 10k continuations of three inspected 100k parents, so its effects are conditional on those parents. Each comparison retains three paired seeds. The native capped base graph and self-messages are preserved; architecture and protocol details appear in Appendix~\ref{sec:implementation-details}.

Goop shows that training support matters even for an allocator with no learned ranking. Under the same random25 evaluation policy, full-rollout MSE falls from $0.301\pm0.169$ to $0.129\pm0.031$ after graph exposure, improving in every seed (Figure~\ref{fig:goop-exposure-placement}A). This establishes a benefit of exposure for a fixed expansion policy. It does not yet show that residual scores can select useful edges, or that expansion improves on the native base graph.

% INSERT unchanged Figure~\ref{fig:goop-exposure-placement} here, preferably
% with a shorter thesis-led caption; move full-rollout guard counts to that test.

\subsection{Does residual risk place the budget well?}
To isolate placement, we evaluate each policy on the same 150 observed test histories per model. Random25, speed25, risk25 and RMS25 retain the same quarter-annulus budget at radius ratio 1.267; base and dense control whether to expand at all. Speed uses displacement magnitude, RMS uses relative velocity over native neighbors, and risk uses scores from the preceding observed base history. These comparisons hold the input state and model fixed.

The exposure gain does not resolve Goop's placement problem. Previous-observed risk remains worse than random in every mixed-model seed, although its mean deficit narrows from $(6.39\pm1.45)\times10^{-10}$ to $(2.84\pm3.62)\times10^{-10}$ (Figure~\ref{fig:goop-exposure-placement}B; Table~\ref{tab:goop-exposure-placement}). Its correlation with base error is $.214\pm.088$, but its correlation with its own sparse-action benefit is $.001\pm.048$. The score identifies difficult predictions much better than it identifies predictions that its chosen expansion improves.

WaterDrop tests whether exposure can change the placement ordering. It reverses the observed-test risk-minus-random ordering in every seed: the mean gap changes from $(+1.62\pm1.07)\times10^{-10}$ to $(-0.99\pm0.67)\times10^{-10}$. The validation gap also improves, but mixed risk remains worse than random there on average. Sand shows a third pattern: exposure narrows its test deficit in every seed, from $(+4.76\pm3.80)\times10^{-9}$ to $(+1.87\pm1.04)\times10^{-9}$, while risk continues to lose to random in both arms. A favorable training interaction can therefore coexist with ineffective placement.

% INSERT unchanged Table~\ref{tab:goop-exposure-placement} here.

Earlier WaterDrop controls help explain this distinction. We measure action benefit as base-graph loss minus expanded-policy loss on the same history and target. At 100k updates, faithful and NLL risk correlate with base error ($.357\pm.060$ and $.411\pm.023$), but weakly with realized sparse-action benefit ($-.021\pm.049$ and $-.052\pm.013$); random has lower one-step error in every seed. Restoring trained self-messages lowers error without reversing that ordering. The corresponding action decomposition shows why improved alignment with the target is insufficient: risk's larger squared prediction change outweighs its alignment gain in every full-model seed (Appendices~\ref{sec:actual-action-benefit} and~\ref{sec:graph-bridge}). Predicting a residual and predicting how an intervention changes that residual are distinct tasks.

\subsection{Do the placement findings survive feedback?}
An observed-history comparison removes a major source of variation, but a deployed simulator repeatedly predicts its own inputs. Its cached-risk controller also uses scores from its preceding selected graph, whereas the observed test scores a preceding base graph. We therefore evaluate all declared policies over each complete forecast horizon: H395 for Goop, H995 for WaterDrop and H314 for Sand. Frames average within trajectories, trajectories within seeds, and reported means and sample SDs retain all three seeds.

WaterDrop illustrates why the observed ordering alone is insufficient. All 810 required rollout outcomes complete, and exposure improves random, dense, speed and cached risk in every seed. Yet the autonomous risk-minus-random training interaction is $+.00284\pm.00466$, with mixed signs across seeds. The observed-test reversal does not yield a consistent autonomous allocation gain.

Sand is the stronger counterexample. All 1,080 outcomes complete, but exposure worsens cached-risk error in every seed ($+.00917\pm.01201$) and widens its risk-minus-random gap in every seed ($+.01975\pm.03062$). This happens despite the smaller observed-history deficit. Random improves in only two seeds, and native base beats random in every seed of both arms. Dense exposure improves every seed, but dense still loses to both native base and random. These comparisons separate a model's response to training from the value of expanding its graph.

Goop's full-horizon risk conclusion is limited by three candidate-pair guard failures in mixed seed 2, one each under base, dense and cached risk. The affected mixed-arm means and the cached-risk training interaction remain undefined. The two defined risk-minus-random seed comparisons have opposite signs. Among complete mixed-arm policies, speed25 and RMS25 beat random25 in every seed, with MSE $.112\pm.027$ and $.125\pm.032$, respectively (Table~\ref{tab:goop-rollouts}).

% INSERT unchanged Table~\ref{tab:goop-rollouts} here. Preserve every policy,
% undefined effect and unsuccessful-outcome count in its current payload.

Lower rollout MSE also leaves substantial physical mismatch. Goop's mixed random, speed and RMS policies place 19.55--22.16\% of particles more than $10^{-6}$ outside the metadata box, compared with $0.0632\%$ for truth. These geometric diagnostics do not test conservation. Complete policy results, excursions and failure accounting are reported in Appendices~\ref{sec:goop-graph-exposure}, \ref{sec:waterdrop-110k-exposure} and~\ref{sec:sand-graph-exposure}; Figure~\ref{fig:qualitative-goop2d} shows a fixed metadata-selected example.

\subsection{What computation does adaptation require?}
Caching reuses the preceding simulator output, avoiding a new scoring pass at each subsequent forecast. It still incurs candidate construction, selection and additional message passing. The observed-history risk test includes its separate scoring pass, so its timing measures a different computation. A fixed optional-pair budget compares placement on a common state; full rollouts instead develop different geometries, which changes even the mandatory graph. We report measured costs with the policy results, but neither rollout edge counts nor fixed-order timings establish a wall-clock advantage. The follow-up designs are exploratory and use three seeds; their different training interventions and horizons do not isolate material effects.
